\section{Expert-guided graph ensembles and their intervention limits}
\label{exec20:expert-scope}

A bank building a diagram from several discovery algorithms will often face
agreement about a relationship but disagreement about its direction. The
methods may also miss different relationships because they use different
scores, tests and structural restrictions. Asking a domain expert to resolve
selected disagreements is a plausible use of an assistant. The design question
is what the expert is being asked, how its answer changes the graph, and whether
the resulting graph supports the bank's intervention query.

Tench and Demeester propose a dynamic expert-guided ensemble that distinguishes
adjacency from orientation. It counts how many component graphs contain any
edge between a pair, considers pairs in descending frequency, and adds directed
edges subject to an acyclicity check. A lower adjacency threshold determines
which pairs are eligible. For eligible pairs without a strict majority, it asks
the expert whether a plausible causal relationship exists. A separate threshold
controls when direction votes suffice and when expert orientation is requested.
The procedure uses equal component weights and a finite query budget. This is
greedy graph construction; the term model averaging does not establish a
Bayesian posterior over causal models. \citep{exec20_expert_averaging}

\subsection{An expert's causal ordering is not a direct edge}
\label{exec20:expert-semantics}

The source deliberately permits an expert to answer about a direct or indirect
causal path. This is weaker than knowing whether two nodes are adjacent in the
full graph. In a stipulated offer--uptake--balance chain
$O\rightarrow U\rightarrow Y$, changing the offer may affect balance through
uptake even though $O\rightarrow Y$ is not a direct edge. An expert can answer
the path question correctly while accepting a spurious direct connection
proposed by a component method. The paper recognizes this distinction in its
simulated expert and compares it with an expert that knows direct edges.
\citep{exec20_expert_averaging,exec20_expert_averaging_reference}

The proposed bank interface should therefore carry the proposition being
elicited: adjacency, ancestry, orientation conditional on adjacency, or a
temporal exclusion. It should also permit insufficient knowledge. A reply about
plausible ancestry cannot silently become proof of a direct edge, and a forced
binary answer cannot represent every legitimate uncertainty state. The set of
variables matters: marginalizing uptake changes the observed-variable graph
even when the underlying mechanism is unchanged.

The method's output is a DAG. Its component comparison includes FCI, while its
noisy benchmarks can contain latent confounding. A directed output cannot by
itself retain every bidirected or unresolved endpoint in a latent graph. Before
an adapter is proposed, the conversion from each component's graph type must be
explicit, with unsupported endpoints preserved or refused. Acyclicity checks
alone do not establish compatibility with an input equivalence class or identify
an intervention in the presence of hidden common causes.
\citep{exec20_expert_averaging}

\subsection{Frequency, diversity and structural compatibility}
\label{exec20:expert-votes}

For one binary question, suppose eleven voters each have probability $3/4$ of
being correct. Under independent errors their majority is correct with
probability
\begin{equation}
 \sum_{k=6}^{11}\binom{11}{k}
       \left(\frac34\right)^k\left(\frac14\right)^{11-k}
 =\frac{1012581}{1048576}\simeq0.966.
 \label{exec20:expert-majority}
\end{equation}
If all eleven instead copy the same $3/4$-correct signal, majority correctness
remains $3/4$. These are two authored probability laws with the same individual
accuracy. They show why individual competence alone cannot justify improvement
toward certainty. The paper itself acknowledges that algorithms share data and
often share statistical principles. Its discussion motivates diversity, but
does not estimate a calibrated graph probability from vote frequency.
\citep{exec20_expert_averaging,exec20_expert_averaging_reference}

Even individually valid DAGs can produce incompatible edge majorities. Consider
the three graphs
\begin{equation}
 \begin{split}
 G_1&=\{A\rightarrow B,\ B\rightarrow C\},\\
 G_2&=\{B\rightarrow C,\ C\rightarrow A\},\\
 G_3&=\{C\rightarrow A,\ A\rightarrow B\}.
 \end{split}
 \label{exec20:expert-cycle}
\end{equation}
Each is acyclic, yet every edge of $A\rightarrow B\rightarrow C\rightarrow A$
has frequency $2/3$. A greedy acyclicity rule must discard or reconsider an
edge. Its result can depend on tie handling and earlier choices. Returning one
DAG makes the output structurally admissible under that rule, but it neither
preserves every component claim nor proves which disputed edge is false.
\citep{exec20_expert_averaging_reference}

\subsection{The reported evidence is informative and conditional}
\label{exec20:expert-evidence}

The main experiments use seven benchmark network structures with synthetic
categorical samples, clean and corrupted versions, four sample sizes from 100
to 100,000, and twenty column orderings. The principal simulated expert is
correct on 70\% of existence questions and 90\% of orientation questions.
Thresholds of zero for adjacency and 0.9 for orientation, with a budget of twice
the node count, were selected using separate randomly generated networks. The
paper reports broad gains in its graph scores, alongside a precision--recall
tradeoff and varying comparison denominators when algorithms fail or return
undefined precision. These are reported experiments; they are not newly
reproduced results in this proposal. \citep{exec20_expert_averaging}

The appendix tests a narrower real-model setting on five clean networks,
using \texttt{gpt-5-nano-}\allowbreak\texttt{2025-08-07} and excluding Pathfinder
because its variable names lack the needed semantics and Formed because many
edges represent diagnostic rather than causal relations. Expert responses are
cached across seeds and sample sizes. Both argument orders are queried, with
disagreement causing a fallback to component voting. The reported per-response
orientation and existence accuracies are not the accuracy of that combined
fallback policy. \citep{exec20_expert_averaging}

\begin{center}
\begin{tabular}{lrr}
\hline
Reported case & Orientation accuracy & Existence accuracy\\
\hline
Mean across five networks & 93.3\% & 66.4\%\\
Property & 83.9\% & 43.9\%\\
Alarm & 82.8\% & 52.2\%\\
\hline
\end{tabular}
\end{center}
Property's existence accuracy falls below the paper's better-than-random expert
assumption. A useful mean therefore cannot replace domain-specific evidence,
and accuracy on all possible pairs need not equal accuracy on the difficult
pairs selected adaptively by an ensemble. Cached responses also do not become
twenty independent expert realizations because the numerical algorithms are
rerun twenty times. The appendix's near reconstruction of the Asia network
raises a memorization concern; unsuccessful recall of another full network
cannot establish that its concepts or fragments were absent from training.
\citep{exec20_expert_averaging}

The source states that it supplies no formal graph-recovery guarantee and has
not empirically qualified continuous or mixed observations. Its cost discussion
includes component discovery, expert requests and structural-distance scoring.
The proposed comparison should likewise retain timeout, memory and invalid-
output denominators, distinguishing a conditional score on successful outputs
from all-request performance. The source's proposed imputation sensitivity
analysis does not supply the missing runtime results or a bank operating cost.

\subsection{Resolve the printed query contract before reuse}
\label{exec20:expert-budget}

The prose says that a fixed budget bounds queries and exhausted decisions fall
back to majority voting. The printed Algorithm~1 leaves two details unresolved.
For an existence query, the counter increment lies inside the branch reached
only when the expert accepts the connection. With three eligible minority
connections, budget one, and three rejections, a literal trace makes three
expert calls while the printed counter remains zero. Separately, when an
accepted adjacency has two legal directions, neither reaches the orientation
threshold, and no budget remains, the printed conditional chain has no final
branch adding the majority direction. An example with direction votes
$3/4$ and $1/4$, threshold $9/10$ and budget zero therefore leaves the edge
unassigned under the printed chain. \citep{exec20_expert_averaging_reference}

These are specification discrepancies, not verified defects in the released
implementation. Resolving them requires an immutable source version and its
tests and call logs. A proposed adapter should count every attempted external
call, distinguish logical questions from the two argument-order requests,
account for retries and invalid answers, and define the behavior at exhaustion.
Caching needs a key that binds the variable meanings, full request, model
snapshot and relevant settings. Otherwise a cost-saving cache can silently
reuse evidence for a different scientific question. The method and data lead
in the paper are useful acquisition routes; package names and reported license
labels are not sufficient permission or provenance for direct adoption.

\subsection{A selected graph still has to answer the intervention query}
\label{exec20:expert-identification}

For a separate exact example, consider two standardized observational summaries
$X,Y$ whose temporal ordering has not been established. Let
$\rho=3/5$. The following two zero-mean Gaussian structural models use
independent errors within each model:
\begin{equation}
 \begin{aligned}
 \mathcal M_1:\quad &X=\epsilon_X,\quad Y=\rho X+\epsilon_Y,\\
 &\operatorname{Var}(\epsilon_X)=1,\quad
 \operatorname{Var}(\epsilon_Y)=1-\rho^2,\\[2pt]
 \mathcal M_2:\quad &Y=\eta_Y,\quad X=\rho Y+\eta_X,\\
 &\operatorname{Var}(\eta_Y)=1,\quad
 \operatorname{Var}(\eta_X)=1-\rho^2.
 \end{aligned}
 \label{exec20:expert-worlds}
\end{equation}
Both imply covariance
$\left(\begin{smallmatrix}1&3/5\\3/5&1\end{smallmatrix}\right)$ and hence the
same observational distribution. Their effects on the mean of $Y$ of increasing
the intervention level of $X$ by one unit differ:
\begin{equation}
 \tau(\mathcal M_1)=\frac35,\qquad \tau(\mathcal M_2)=0.
 \label{exec20:expert-effects}
\end{equation}
More algorithms fitted to the same distribution do not identify the missing
direction. Truthful temporal restrictions, suitable experiments or independently
justified structural knowledge can eliminate a candidate; that extra information
must be stated. The example does not override a known banking information clock.
\citep{exec20_expert_averaging_reference}

Suppose a one-unit increase in $Y$ is provisionally valued at USD~100 and the
intervention costs USD~40. The two compatible margins are USD~20 and USD~$-40$.
Treating a 90\% graph vote as probability for $\mathcal M_1$ gives effect
$27/50$ and margin USD~14, but the vote has supplied no probabilistic calibration
that warrants that calculation as expected value. A decision under unresolved
model ambiguity must retain the compatible effects or adopt an explicitly
justified probability model. This is an authored planning illustration, not
evidence of model performance or bank returns. \citep{exec20_expert_averaging_reference}

The proposal should therefore retain expert-guided averaging as a candidate
for generating and reviewing diagrams. A public prototype comparison can use
known generated mechanisms, withheld domains, correlated component failures,
wrong expert assertions, budget exhaustion and incompatible graph endpoints.
It should compare static aggregation, selective expert queries, a numerical
method without expert input and a workflow that preserves unresolved graph
classes. The outcomes should include the correctness of the intervention
answer, supported refusal, retained ambiguity, all physical requests and total
review cost. Graph scores remain useful diagnostics, but qualification of an
intervention service requires its own evidence.

